% !TeX root = ../高等代数教材（谢启鸿）笔记.tex
% 由主文档通过 \input 载入；请编译主文档。
\setcounter{chapter}{3} % 保留原书章号。
\chapter{线性映射}
\label{chap:linear-maps}

\section{线性映射的概念}
\label{sec:linear-map-concept}

% 依据教材第四章第 4.1--4.5 节（原书第 183--207 页）整理；不收录章末复习题。
线性空间描述可以相加、可以数乘的对象，\emph{线性映射}则描述保持这两种运算的对应。本章先在不选坐标的情形下研究这种对应，再选取基，将它转化为矩阵运算；像、核与不变子空间将说明矩阵的秩和分块形式背后的空间结构。以下各线性空间及其线性映射均在同一数域 $K$ 上讨论；涉及维数公式或表示矩阵时，再明确有限维假设。

先回顾一般映射的语言。映射 $f:X\to Y$ 要求每个 $x\in X$ 都有唯一确定的像 $f(x)\in Y$；$X$ 是定义域，$Y$ 是陪域，像集为 $f(X)=\{f(x):x\in X\}\subseteq Y$。如果 $f(x)=f(x')$ 总能推出 $x=x'$，称 $f$ 为\emph{单射}（injection）；如果 $f(X)=Y$，称为\emph{满射}（surjection）；既单又满的映射称为\emph{双射}（bijection），也称一一对应。单射保证每个元素至多有一个原像，满射保证陪域的每个元素至少有一个原像，双射则保证恰有一个原像。映射是否满射依赖于陪域，例如 $x\mapsto x^2$ 作为 $\mathbb R\to\mathbb R$ 的映射既不单也不满，作为 $\mathbb R\to[0,+\infty)$ 的映射则满而不单。

若 $f:X\to Y$、$g:Y\to Z$，定义\emph{复合映射}（composition）$g\circ f:X\to Z$ 为 $(g\circ f)(x)=g(f(x))$，即先作用 $f$，再作用 $g$。复合要求前一个映射的像落在后一个映射的定义域中；在相应定义域、陪域匹配时，复合满足结合律，因为
\[
    ((h\circ g)\circ f)(x)=h(g(f(x)))
    =(h\circ(g\circ f))(x).
\]
因此可以略去复合的括号，必要时也将 $g\circ f$ 简记为 $gf$。$X$ 到自身的映射称为变换，其中 $I_X(x)=x$ 称为\emph{恒等映射}（identity map）。定义域、陪域相同的两个映射相等，是指它们在每个元素上的值都相等。

\begin{proposition}
\label{prop:map-two-sided-inverse}
映射 $f:X\to Y$ 是双射，当且仅当存在 $g:Y\to X$，使
\[
    g\circ f=I_X,\qquad f\circ g=I_Y.
\]
此时 $g$ 唯一，称为 $f$ 的\emph{逆映射}（inverse map），记作 $f^{-1}$。
\end{proposition}

\begin{proof}
若 $f$ 是双射，将 $g(y)$ 定义为 $y$ 的唯一原像，就有这两个等式。反之，若两个等式成立，则 $f(x)=f(x')$ 蕴含 $g(f(x))=g(f(x'))$，从而 $x=x'$；对任意 $y\in Y$，又有 $y=f(g(y))$，故 $f$ 也满射。如果 $g,h$ 都满足这两个等式，则
\[
    g=g\circ(f\circ h)=(g\circ f)\circ h=h.\qedhere
\]
\end{proof}

上述证明还表明，只有 $g\circ f=I_X$ 时即可推出 $f$ 单射，只有 $f\circ g=I_Y$ 时即可推出 $f$ 满射；一般不能由单侧等式断言 $f$ 可逆。

\begin{definition}
\label{def:linear-map}
设 $V,U$ 是数域 $K$ 上的线性空间。映射 $T:V\to U$ 若满足
\[
    T(u+v)=T(u)+T(v),\qquad T(av)=aT(v)
    \quad(u,v\in V,\ a\in K),
\]
称为\emph{线性映射}（linear map）。$V$ 到自身的线性映射称为\emph{线性变换}（linear transformation）或线性算子；双射的线性映射称为\emph{线性同构}（linear isomorphism），$V$ 到自身的线性同构称为\emph{自同构}（automorphism）。
\end{definition}

\begin{proposition}
\label{prop:linear-map-basic}
$T:V\to U$ 为线性映射的等价条件是
\begin{equation}
\label{eq:linear-map-two-term}
    T(au+bv)=aT(u)+bT(v)
    \quad(u,v\in V,\ a,b\in K).
\end{equation}
线性映射满足 $T(0_V)=0_U$、$T(-v)=-T(v)$，并保持任意有限线性组合：
\[
    T\left(\sum_{i=1}^{s}a_iv_i\right)
    =\sum_{i=1}^{s}a_iT(v_i).
\]
线性同构的逆映射仍是同一数域上的线性映射。
\end{proposition}

\begin{proof}
由加法性和齐次性直接得到式~\eqref{eq:linear-map-two-term}；反之，分别取 $a=b=1$ 以及 $b=0$，即可得到定义中的两个条件。由 $T(0\cdot v)=0\cdot T(v)$ 可得 $T(0_V)=0_U$，取标量 $-1$ 得 $T(-v)=-T(v)$，对项数归纳即得有限线性组合的公式。若 $T$ 是双射，对 $y,z\in U$，令 $u=T^{-1}(y)$、$v=T^{-1}(z)$，则
\[
    T(au+bv)=ay+bz,
    \qquad T^{-1}(ay+bz)=aT^{-1}(y)+bT^{-1}(z),
\]
故 $T^{-1}$ 线性。
\end{proof}

\begin{example}
零映射 $0_{V,U}(v)=0_U$、恒等变换 $I_V(v)=v$、固定 $c\in K$ 后的\emph{纯量变换}（scalar transformation）$v\mapsto cv$ 都是线性的。给定 $A\in M_{m\times n}(K)$，映射
\[
    L_A:K_n\to K_m,\qquad x\mapsto Ax
\]
也是线性的，这里 $K_n,K_m$ 均采用列向量。将行向量 $(a_1,\ldots,a_n)$ 对应到列向量 $(a_1,\ldots,a_n)^{\mathsf T}$ 则给出行、列向量空间之间的线性同构。
\end{example}

\begin{solve}
前三种映射的线性性直接由空间的运算律得到；矩阵映射满足 $A(ax+by)=aAx+bAy$；行、列向量的对应保持各分量的加法和数乘，并且转置操作给出其逆映射。这些例子说明，线性性要求保持运算，并不要求映射单射或满射；例如 $V\ne\{0\}$ 时零映射就不是单射。
\end{solve}

\begin{example}
在实线性空间 $C^\infty([0,1])$ 上，求导 $D(f)=f'$ 是线性变换。在 $K_n$ 上，固定 $A\in M_n(K)$ 和 $b\in K_n$，映射 $F(x)=Ax+b$ 为线性变换，当且仅当 $b=0$。
\end{example}

\begin{solve}
求导满足 $(af+bg)'=af'+bg'$，且无限次可微函数的导数仍在原空间中。对第二个映射，若 $F$ 线性，则 $b=F(0)=0$；若 $b=0$，则 $F=L_A$ 线性。特别地，非零平移不是线性变换，不能因为坐标表达式只含一次项就忽略常数项。
\end{solve}

\begin{notealert}[type=caution,title={先明确数域与空间}]
同一个映射的线性性必须相对于指定数域判断。例如复共轭 $z\mapsto\overline z$ 是 $\mathbb C$ 作为实线性空间上的线性变换，却不是复线性变换，因为 $\overline{i\cdot1}=-i\ne i\overline1$。求导是否为某个空间上的变换，也要检查导函数是否仍属于该空间；仅写出运算公式还不足以确定线性变换。
\end{notealert}

\section{线性映射的运算}
\label{sec:linear-map-operations}

矩阵可以相加和数乘，与之对应的线性映射也可以逐点相加和数乘。对于定义域、陪域都相同的线性映射 $S,T:V\to U$ 及 $a\in K$，规定
\begin{equation}
\label{eq:linear-map-add-scale}
    (S+T)(v)=S(v)+T(v),\qquad (aT)(v)=aT(v).
\end{equation}
记 $\mathcal L(V,U)$ 为所有 $V\to U$ 的 $K$-线性映射组成的集合，记 $\mathcal L(V)=\mathcal L(V,V)$。这里映射的零元是零映射，$T$ 的加法逆元为 $(-T)(v)=-T(v)$，均应与空间中的零向量相区别。

\begin{proposition}
\label{prop:linear-map-space}
在式~\eqref{eq:linear-map-add-scale} 的运算下，$\mathcal L(V,U)$ 是 $K$ 上的线性空间。特别地，$V\to K$ 的线性映射称为\emph{线性函数}或\emph{线性泛函}（linear functional），它们组成的空间记为 $V^*=\mathcal L(V,K)$，称为 $V$ 的\emph{代数对偶空间}（algebraic dual space）。
\end{proposition}

\begin{proof}
先验证封闭性。对 $u,v\in V$、$a,b\in K$，
\[
    (S+T)(au+bv)
    =a(S(u)+T(u))+b(S(v)+T(v))
    =a(S+T)(u)+b(S+T)(v),
\]
故 $S+T$ 线性；同理由 $(cT)(au+bv)=a(cT)(u)+b(cT)(v)$ 得数乘封闭。其余运算律逐点归结为 $U$ 中的运算律，例如
\[
    ((R+S)+T)(v)=R(v)+S(v)+T(v)=(R+(S+T))(v).
\]
同样可验证加法交换律、零元与逆元，以及 $1T=T$、$(ab)T=a(bT)$、$a(S+T)=aS+aT$、$(a+b)T=aT+bT$，所以线性空间的全部公理成立。
\end{proof}

教材在一般情形下也称 $V^*$ 为“共轭空间”，在有限维情形称其为对偶空间；这里统一使用代数对偶的名称，它只涉及 $K$-线性函数，并不要求引入内积、连续性或复共轭运算。

若 $T:V\to U$、$S:U\to W$ 都线性，则它们的复合也线性，因为
\[
    (S\circ T)(au+bv)
    =S(aT(u)+bT(v))
    =a(S\circ T)(u)+b(S\circ T)(v).
\]
特别地，在 $\mathcal L(V)$ 中可以将复合作为乘法，记 $ST=S\circ T$。这里乘法次序仍是\emph[caution]{先作用右侧的变换，再作用左侧的变换}，不能按书写顺序从左向右作用。

\begin{definition}
\label{def:associative-unital-algebra}
$K$ 上的线性空间 $\mathcal A$ 若还具有乘法，满足结合律，存在双侧恒等元 $1_{\mathcal A}$，且对 $a,b,c\in\mathcal A$、$\lambda\in K$ 有
\[
    a(b+c)=ab+ac,\qquad (a+b)c=ac+bc,
    \qquad (\lambda a)b=\lambda(ab)=a(\lambda b),
\]
则称为 $K$ 上的\emph{含幺结合代数}（unital associative algebra）。本章依教材将这种结构简称为代数，不要求乘法交换。
\end{definition}

\begin{theorem}
\label{thm:endomorphism-algebra}
以复合为乘法，$\mathcal L(V)$ 是 $K$ 上的含幺结合代数，其乘法恒等元为 $I_V$。
\end{theorem}

\begin{proof}
线性空间结构已由命题~\ref{prop:linear-map-space} 建立，乘法封闭性由复合的线性性得到。结合律是一般映射的复合结合律，$I_VT=TI_V=T$。分配律逐点验证：
\[
    (R(S+T))(v)=R(S(v)+T(v))=(RS)(v)+(RT)(v),
\]
\[
    ((R+S)T)(v)=R(T(v))+S(T(v))=(RT)(v)+(ST)(v).
\]
前一个等式使用 $R$ 的加法性。类似地，$((\lambda S)T)(v)=\lambda S(T(v))=(S(\lambda T))(v)$，即得数乘相容性。
\end{proof}

\begin{proposition}
\label{prop:linear-map-powers}
对任意 $T\in\mathcal L(V)$，规定 $T^0=I_V$，$T^n$（$n\geq1$）为 $n$ 个 $T$ 的复合，则对非负整数 $r,s$ 有
\[
    T^rT^s=T^{r+s},\qquad (T^r)^s=T^{rs}.
\]
若 $T$ 可逆，再规定 $T^{-n}=(T^{-1})^n$，上述公式便对所有整数指数成立，且 $T^{-n}=(T^n)^{-1}$。若 $S,T$ 都可逆，$c\in K\setminus\{0\}$，则
\[
    (ST)^{-1}=T^{-1}S^{-1},\qquad (cT)^{-1}=c^{-1}T^{-1}.
\]
\end{proposition}

\begin{proof}
非负整数幂的公式由结合律得到；逆幂的公式再使用 $TT^{-1}=T^{-1}T=I_V$ 消去互逆因子。对复合的逆，直接计算
\[
    (ST)(T^{-1}S^{-1})=I_V,
    \qquad (T^{-1}S^{-1})(ST)=I_V.
\]
命题~\ref{prop:map-two-sided-inverse} 给出所需逆映射；数乘情形同理由两侧相乘验证。
\end{proof}

\begin{example}
\label{ex:linear-map-noncommuting}
在 $K^2$ 上定义 $S(x,y)=(y,0)$、$T(x,y)=(0,x)$。则
\[
    ST(x,y)=(x,0),\qquad TS(x,y)=(0,y),
\]
故 $ST\ne TS$。又有 $S^2=T^2=0$、$(ST)^2=ST\ne0$，所以一般不能把 $(ST)^2$ 写成 $S^2T^2$。
\end{example}

只有额外知道 $ST=TS$ 时，才能通过交换中间因子得到 $(ST)^n=S^nT^n$，以及通常的二项式展开。在线性空间 $\mathcal L(V)$ 中建立加法和数乘，并不会自动使其乘法具有数的交换律；负指数也只有在变换可逆时才有定义。

\section{线性映射与矩阵}
\label{sec:linear-map-matrix}

要将抽象映射化为可计算的数据，首先需要知道哪些数据足以唯一确定一个线性映射。线性性保证：只要给出定义域一组基的像，其他向量的像就随其线性组合确定。

\begin{lemma}
\label{lem:linear-map-basis-images}
设 $V$ 的一组基为 $E=(e_1,\ldots,e_n)$，$U$ 是同一数域上的线性空间。给定任意 $u_1,\ldots,u_n\in U$，存在唯一线性映射 $T:V\to U$，使 $T(e_j)=u_j$（$1\leq j\leq n$）。因此，两个线性映射相等，当且仅当它们在这组基的各向量上取值相同。
\end{lemma}

\begin{proof}
每个 $v\in V$ 唯一表示为 $v=\sum_{j=1}^n x_je_j$，定义 $T(v)=\sum_{j=1}^n x_ju_j$，就得到一个定义良好的映射。对向量的线性组合使用坐标的加法、数乘，可验证此映射线性；取 $v=e_j$ 即得规定的像。反之，任何满足要求的线性映射都必须把 $\sum_jx_je_j$ 映为 $\sum_jx_ju_j$，故唯一。
\end{proof}

\begin{definition}
\label{def:linear-map-representation}
设 $\dim V=n$、$\dim U=m$，分别固定有序基 $E=(e_1,\ldots,e_n)$、$F=(f_1,\ldots,f_m)$。若 $T:V\to U$ 线性，写成
\[
    T(e_j)=\sum_{i=1}^m a_{ij}f_i\quad(1\leq j\leq n),
\]
则矩阵 $A=(a_{ij})_{m\times n}$ 称为 $T$ 在基 $E,F$ 下的\emph{表示矩阵}（representation matrix），记作 $[T]_{F\leftarrow E}$。也就是说，\emph{第 $j$ 列是 $T(e_j)$ 在陪域基 $F$ 下的坐标列向量}。
\end{definition}

教材将 $T(e_j)$ 中 $f_i$ 的系数记为 $a_{ji}$，因而需要对其逐行列出的系数阵转置；这里直接把该系数记为 $a_{ij}$，使第一下标表示行、第二下标表示列。两种记法给出的表示矩阵相同。以下统一采用坐标列向量：若 $v=\sum_jx_je_j$，记 $[v]_E=(x_1,\ldots,x_n)^{\mathsf T}$，则
\begin{equation}
\label{eq:linear-map-coordinate}
    [T(v)]_F=A[v]_E.
\end{equation}
事实上，$T(v)=\sum_jx_j\sum_ia_{ij}f_i=\sum_i(\sum_ja_{ij}x_j)f_i$，恰好给出矩阵乘法。反过来，式~\eqref{eq:linear-map-coordinate} 对所有 $v$ 成立时，逐一取 $v=e_j$，就恢复了表示矩阵的各列。

\begin{theorem}
\label{thm:linear-map-matrix-isomorphism}
固定上述基后，映射
\[
    \Phi_{F,E}:\mathcal L(V,U)\longrightarrow M_{m\times n}(K),
    \qquad T\longmapsto[T]_{F\leftarrow E}
\]
是线性同构。因此 $\dim\mathcal L(V,U)=mn$。若 $\eta_E(v)=[v]_E$、$\eta_F(u)=[u]_F$ 是坐标同构，则
\[
    \eta_F\circ T=L_A\circ\eta_E,
    \qquad T=\eta_F^{-1}\circ L_A\circ\eta_E.
\]
\end{theorem}

\begin{proof}
引理~\ref{lem:linear-map-basis-images} 表明，每个矩阵的各列作为基向量的像坐标，唯一确定一个线性映射，所以 $\Phi_{F,E}$ 是双射。由于 $(S+T)(e_j)=S(e_j)+T(e_j)$、$(cT)(e_j)=cT(e_j)$，取各列的坐标可得
\[
    [S+T]_{F\leftarrow E}=[S]_{F\leftarrow E}+[T]_{F\leftarrow E},
    \qquad [cT]_{F\leftarrow E}=c[T]_{F\leftarrow E}.
\]
故 $\Phi_{F,E}$ 线性，维数结论由矩阵空间的标准基得到。最后两个等式就是式~\eqref{eq:linear-map-coordinate} 的映射写法。
\end{proof}

\begin{theorem}
\label{thm:linear-map-composition-matrix}
设 $T:V\to U$、$S:U\to W$ 线性，$E,F,G$ 分别是三个有限维空间的有序基，且
\[
    A=[T]_{F\leftarrow E},\qquad B=[S]_{G\leftarrow F}.
\]
则
\begin{equation}
\label{eq:linear-map-composition-matrix}
    [S\circ T]_{G\leftarrow E}=BA.
\end{equation}
特别地，固定 $V$ 的同一组基用于定义域和陪域后，$T\mapsto[T]_E$ 是 $\mathcal L(V)$ 与 $M_n(K)$ 之间保持乘法及恒等元的代数同构；$T$ 可逆当且仅当 $[T]_E$ 可逆，此时
\[
    [I_V]_E=I_n,\qquad [T^{-1}]_E=[T]_E^{-1}.
\]
\end{theorem}

\begin{proof}
对任意 $v\in V$，
\[
    [S(T(v))]_G=B[T(v)]_F=BA[v]_E,
\]
故复合的表示矩阵为 $BA$。恒等映射将每个基向量保持不变，其矩阵为 $I_n$；若 $T$ 可逆，将 $TT^{-1}=T^{-1}T=I_V$ 取矩阵，得到逆矩阵公式。反之，若 $A=[T]_E$ 可逆，用引理~\ref{lem:linear-map-basis-images} 取表示矩阵为 $A^{-1}$ 的线性变换 $S$，则 $ST,TS$ 的矩阵均为 $I_n$，从而 $ST=TS=I_V$，故 $T$ 可逆。
\end{proof}

式~\eqref{eq:linear-map-composition-matrix} 给出了矩阵乘法的空间意义。中间空间 $U$ 的基必须在两个矩阵中采用同一个 $F$，这样 $T$ 的输出坐标才能直接成为 $S$ 的输入坐标；只有矩阵尺寸相配还不足以替代这一约定。

\begin{example}
\label{ex:linear-map-derivative-matrix}
设 $V=K[x]_{<n}$ 是次数小于 $n$ 的多项式与零多项式组成的空间，$n\geq2$，在基 $E=(1,x,\ldots,x^{n-1})$ 下求导变换 $D$ 的矩阵。
\end{example}

\begin{solve}
由 $D(1)=0$、$D(x^j)=jx^{j-1}$（$1\leq j\leq n-1$），第 $j+1$ 列只有第 $j$ 行的元素 $j$ 非零。因此 $[D]_E=(d_{ij})$，其中 $d_{j,j+1}=j$，其余元素为零；例如 $n=4$ 时，
\[
    [D]_E=\begin{pmatrix}
        0&1&0&0\\
        0&0&2&0\\
        0&0&0&3\\
        0&0&0&0
    \end{pmatrix}.
\]
这个矩阵编码了求导降低次数的作用，也给出 $D^n=0$，因为任意次数小于 $n$ 的多项式求导 $n$ 次都为零。
\end{solve}

表示矩阵依赖于基，线性映射本身不依赖于基。为写清换基方向，设定义域和陪域的新基分别为 $E'=(e'_1,\ldots,e'_n)$、$F'=(f'_1,\ldots,f'_m)$，并约定
\[
    E'=EP,\qquad F'=FQ,
\]
其中 $P,Q$ 的各列分别是新基向量在相应旧基下的坐标。因此
\[
    [v]_E=P[v]_{E'},\qquad [u]_F=Q[u]_{F'}.
\]
若 $A=[T]_{F\leftarrow E}$、$B=[T]_{F'\leftarrow E'}$，则对任意 $v$，
\[
    QB[v]_{E'}=[T(v)]_F=AP[v]_{E'}.
\]
由于 $[v]_{E'}$ 遍历全部 $K_n$，取标准单位列向量即可得到 $QB=AP$，从而
\begin{equation}
\label{eq:linear-map-general-basis-change}
    B=Q^{-1}AP.
\end{equation}
这同时说明，换基公式中的左右因子分别负责输出坐标和输入坐标，不能只凭记忆猜测逆矩阵的位置。

\begin{definition}
\label{def:matrix-similarity}
同一数域上的 $n$ 阶方阵 $A,B$，若存在可逆矩阵 $P$ 使
\begin{equation}
\label{eq:linear-map-similarity}
    B=P^{-1}AP,
\end{equation}
则称 $A,B$ \emph{相似}（similar），记为 $A\sim B$。
\end{definition}

\begin{theorem}
\label{thm:linear-map-similarity}
有限维空间上同一线性变换在不同基下的表示矩阵相似；反之，两个相似矩阵可以表示同一线性变换在两组基下的作用。相似是一种等价关系。
\end{theorem}

\begin{proof}
对线性变换，定义域和陪域都取旧基 $E$、新基 $E'=EP$，在式~\eqref{eq:linear-map-general-basis-change} 中令 $Q=P$，即得式~\eqref{eq:linear-map-similarity}。反之，给定 $B=P^{-1}AP$，在 $K_n$ 的标准基 $E$ 下取 $T=L_A$，再取新基 $E'=EP$；$P$ 可逆保证这确为一组基，且 $[T]_{E'}=B$。

自反性由 $A=I_n^{-1}AI_n$ 得到；若 $B=P^{-1}AP$，则 $A=(P^{-1})^{-1}BP^{-1}$，故关系对称；若另有 $C=Q^{-1}BQ$，则 $C=(PQ)^{-1}A(PQ)$，故关系传递。
\end{proof}

相似矩阵有相同的秩、行列式和迹。秩不变是因为左右乘可逆矩阵不改变秩；行列式满足 $\det(P^{-1}AP)=\det A$；迹利用 $\operatorname{tr}(XY)=\operatorname{tr}(YX)$ 得
\[
    \operatorname{tr}(P^{-1}AP)=\operatorname{tr}(APP^{-1})=\operatorname{tr}A,
\]
其中迹的交换公式可由 $\sum_i\sum_jx_{ij}y_{ji}=\sum_j\sum_iy_{ji}x_{ij}$ 核实。因而这些量可以看作线性变换本身的量。对于一般 $V\to U$ 的映射，定义域、陪域允许独立换基，得到的是式~\eqref{eq:linear-map-general-basis-change}；只有同一空间上的变换同时采用同一组新基时，才得到相似关系。

\begin{example}
在实平面标准基下，绕原点逆时针旋转 $\theta$ 角的变换 $R_\theta$ 的矩阵为
\[
    [R_\theta]=\begin{pmatrix}
        \cos\theta&-\sin\theta\\
        \sin\theta&\cos\theta
    \end{pmatrix}.
\]
\end{example}

\begin{solve}
标准基向量 $(1,0)^{\mathsf T}$、$(0,1)^{\mathsf T}$ 的像依次为 $(\cos\theta,\sin\theta)^{\mathsf T}$、$(-\sin\theta,\cos\theta)^{\mathsf T}$，将这两个坐标向量作为列即得矩阵。其行列式为 $1$，故旋转可逆，且 $R_\theta^{-1}=R_{-\theta}$。
\end{solve}

\section{线性映射的像与核}
\label{sec:linear-map-image-kernel}

矩阵表示把映射转化为坐标运算，但仍需用不依赖于坐标的对象描述它保留了多少信息、又把哪些向量压缩为零。像刻画所有可能的输出，核刻画无法从输出中区分的输入差异，两者分别联系矩阵的列空间和齐次方程组的解空间。

\begin{definition}
\label{def:linear-map-image-kernel}
设 $T:V\to U$ 线性，称
\[
    \operatorname{Im}T=T(V)=\{T(v):v\in V\}\subseteq U
\]
为 $T$ 的\emph{像}或\emph{像空间}（image），称
\[
    \ker T=\{v\in V:T(v)=0_U\}\subseteq V
\]
为 $T$ 的\emph{核}或\emph{核空间}（kernel）。像空间的维数称为 $T$ 的\emph{秩}（rank），记为 $r(T)=\dim\operatorname{Im}T$；核空间的维数称为 $T$ 的\emph{零度}（nullity）。
\end{definition}

\begin{proposition}
\label{prop:linear-map-image-kernel-subspace}
$\operatorname{Im}T$ 是 $U$ 的子空间，$\ker T$ 是 $V$ 的子空间，并且
\[
    T\text{ 单射}\iff\ker T=\{0_V\},\qquad
    T\text{ 满射}\iff\operatorname{Im}T=U.
\]
当 $U$ 有限维时，后一个条件又等价于 $r(T)=\dim U$。
\end{proposition}

\begin{proof}
$T(0_V)=0_U$ 保证像和核都含零向量。若 $y=T(v)$、$z=T(w)$，则 $ay+bz=T(av+bw)\in\operatorname{Im}T$；若 $v,w\in\ker T$，则 $T(av+bw)=0_U$，故核也对线性组合封闭。若 $T$ 单射，由 $T(v)=T(0_V)$ 得 $v=0_V$，故核为零。反之，若核为零，则
\[
    T(v)=T(w)\ \Longrightarrow\ T(v-w)=0_U
    \ \Longrightarrow\ v-w=0_V,
\]
故 $T$ 单射。满射的判定就是定义；当 $U$ 有限维时，其子空间与它维数相同当且仅当子空间等于全空间，得到最后的判定。
\end{proof}

\begin{proposition}
\label{prop:linear-map-fibers}
对 $u\in U$，若存在 $v_0\in V$ 使 $T(v_0)=u$，则 $u$ 的全部原像为
\begin{equation}
\label{eq:linear-map-fiber}
    T^{-1}(\{u\})=v_0+\ker T
    =\{v_0+w:w\in\ker T\}.
\end{equation}
如果 $u\notin\operatorname{Im}T$，则其原像集为空。
\end{proposition}

\begin{proof}
$T(v)=u=T(v_0)$ 当且仅当 $T(v-v_0)=0_U$，也就当且仅当 $v-v_0\in\ker T$。
\end{proof}

式~\eqref{eq:linear-map-fiber} 中的 $T^{-1}(\{u\})$ 表示集合的原像，不预先要求 $T$ 可逆；当 $u\ne0_U$ 时，这个集合一般是核的平移，并不是线性子空间。在坐标下，它正是非齐次线性方程组“特解加齐次通解”的空间解释。

\begin{definition}
\label{def:linear-map-restriction}
设 $V'\subseteq V$、$U'\subseteq U$ 是子空间，并且 $T(V')\subseteq U'$。采用相同的对应法则 $v\mapsto T(v)$，得到线性映射 $T':V'\to U'$，称为 $T$ 的\emph{限制}（restriction），通常将定义域的限制记为 $T|_{V'}$，其陪域由上下文指定。
\end{definition}

限制仍线性，因为相应的加法、数乘从原空间继承，而且
\begin{equation}
\label{eq:linear-map-restriction-kernel}
    \ker T'=V'\cap\ker T,\qquad
    \operatorname{Im}T'=T(V').
\end{equation}
因此单射限制后仍单，但满射限制后未必仍满；若将陪域改为 $T(V')$，它又按定义成为满射。例如实平面上的旋转 $R_\theta$ 将 $x$ 轴 $V'$ 映到与 $x$ 轴成 $\theta$ 角的直线 $U'$，其限制 $V'\to U'$ 是线性同构，作为 $V'\to\mathbb R^2$ 的映射则不是满射。

\begin{theorem}[线性映射维数公式]
\label{thm:linear-map-rank-nullity}
设 $V$ 有限维，$T:V\to U$ 线性，则
\begin{equation}
\label{eq:linear-map-rank-nullity}
    \dim\ker T+\dim\operatorname{Im}T=\dim V.
\end{equation}
\end{theorem}

\begin{proof}
设 $\dim V=n$，取 $\ker T$ 的基 $v_1,\ldots,v_k$，将其扩充为 $V$ 的基 $v_1,\ldots,v_k,v_{k+1},\ldots,v_n$。若核为零，前一组取空集。由于前 $k$ 个基向量的像为零，任意向量的像都是 $T(v_{k+1}),\ldots,T(v_n)$ 的线性组合，故这些像生成 $\operatorname{Im}T$。

若 $\sum_{j=k+1}^n c_jT(v_j)=0_U$，则 $\sum_{j=k+1}^n c_jv_j\in\ker T$，于是存在 $d_1,\ldots,d_k$ 使
\[
    \sum_{j=k+1}^n c_jv_j-\sum_{i=1}^k d_iv_i=0_V.
\]
由扩充后的基线性无关，所有 $c_j,d_i$ 均为零。因此这些像还线性无关，构成像空间的一组基，故 $\dim\operatorname{Im}T=n-k$。当 $k=n$ 时像空间为零，同样得到结论。
\end{proof}

这个证明只要求定义域有限维，陪域不必有限维。它还给出一个具体的同构：令 $W=\operatorname{span}\{v_{k+1},\ldots,v_n\}$，则 $V=\ker T\oplus W$，而限制 $T|_W:W\to\operatorname{Im}T$ 是双射。维数公式说的是核与像的维数之和，并不意味着它们在同一个空间里构成直和；即使 $T$ 是 $V$ 上的变换，也未必有 $V=\ker T\oplus\operatorname{Im}T$。

\begin{theorem}
\label{thm:linear-map-rank-matrix}
设 $\dim V=n$、$\dim U=m$，$A=[T]_{F\leftarrow E}$，则
\begin{equation}
\label{eq:linear-map-rank-matrix}
    \dim\operatorname{Im}T=\operatorname{rank}A,
    \qquad \dim\ker T=n-\operatorname{rank}A.
\end{equation}
更具体地，坐标同构 $\eta_E,\eta_F$ 给出
\[
    \eta_E(\ker T)=\{x\in K_n:Ax=0\},
    \qquad \eta_F(\operatorname{Im}T)=\operatorname{Col}(A),
\]
其中 $\operatorname{Col}(A)$ 表示 $A$ 的列空间。
\end{theorem}

\begin{proof}
由式~\eqref{eq:linear-map-coordinate}，$T(v)=0_U$ 等价于 $A[v]_E=0$，且 $[v]_E$ 遍历全部 $K_n$，故第一个集合等式成立。若 $A=(a_1,\ldots,a_n)$ 按列分块，则 $Ax=\sum_jx_ja_j$，全部输出坐标组成的集合恰是这些列向量的张成空间，得到第二个集合等式。将坐标同构限制到相应子空间仍为同构，故 $\dim\operatorname{Im}T=\dim\operatorname{Col}(A)=\operatorname{rank}A$，再由定理~\ref{thm:linear-map-rank-nullity} 得核的维数。
\end{proof}

\begin{corollary}
\label{cor:linear-map-injective-surjective-rank}
在定理~\ref{thm:linear-map-rank-matrix} 的条件下，$T$ 单射当且仅当 $A$ 列满秩，即 $\operatorname{rank}A=n$；$T$ 满射当且仅当 $A$ 行满秩，即 $\operatorname{rank}A=m$。因此单射要求 $n\leq m$，满射要求 $n\geq m$；若 $n=m$，则单射、满射、同构三者等价。
\end{corollary}

\begin{proof}
由命题~\ref{prop:linear-map-image-kernel-subspace}，单射等价于核维数为零，满射等价于像维数为 $m$，再代入式~\eqref{eq:linear-map-rank-matrix} 即可。秩不超过矩阵的行数和列数，给出两个维数限制；当 $n=m$ 时两个满秩条件相同。
\end{proof}

特别地，有限维空间上的线性变换只要单射或满射就已经可逆；这也等价于其任意一组基下的表示矩阵可逆，以及它把一组基映成一组基。后一说法由基向量的像为矩阵各列得到：$n$ 个像构成基，恰好要求表示矩阵的 $n$ 列线性无关。

\begin{example}
\label{ex:linear-map-infinite-dimension}
在无限维空间 $K[x]$ 上，定义求导 $D(f)=f'$ 和零常数项的积分
\[
    J\left(\sum_{j=0}^d a_jx^j\right)
    =\sum_{j=0}^d\frac{a_j}{j+1}x^{j+1}.
\]
则 $D$ 满而不单，$J$ 单而不满，并且
\[
    DJ=I,\qquad JD(f)=f-f(0).
\]
\end{example}

\begin{solve}
$D,J$ 都保持线性组合，直接求导、积分得到上述等式。每个多项式 $f$ 都等于 $D(J(f))$，故 $D$ 满射，但非零常数都在其核中；由 $DJ=I$ 可知 $J$ 单射，而 $J$ 的像中常数项都为零，故常数多项式 $1$ 没有原像。这说明有限维条件不能从单射、满射等价的结论中删除，也再次说明单侧逆不必是双侧逆。
\end{solve}

\begin{example}
\label{ex:linear-map-image-kernel-matrix}
设 $V,U$ 分别为 $K$ 上的五维、四维空间，基分别为 $E=(e_1,\ldots,e_5)$、$F=(f_1,\ldots,f_4)$，线性映射 $T:V\to U$ 的表示矩阵为
\[
    A=\begin{pmatrix}
        1&2&1&-3&2\\
        2&1&1&1&-4\\
        1&1&2&2&-2\\
        2&3&-5&-17&8
    \end{pmatrix}.
\]
求 $\operatorname{Im}T$ 与 $\ker T$ 的基。
\end{example}

\begin{solve}
对 $A$ 作初等行变换，得到行最简形
\[
    R=\begin{pmatrix}
        1&0&0&1&-3\\
        0&1&0&-3&3\\
        0&0&1&2&-1\\
        0&0&0&0&0
    \end{pmatrix}.
\]
主元位于前三列，故 $A$ 的前三列线性无关，且生成整个列空间。因此像空间的一组基为
\[
    T(e_1)=f_1+2f_2+f_3+2f_4,\qquad
    T(e_2)=2f_1+f_2+f_3+3f_4,\qquad
    T(e_3)=f_1+f_2+2f_3-5f_4.
\]
行变换将 $A$ 左乘可逆矩阵，保留各列之间的全部线性关系，所以可以用 $R$ 选择主元列的编号，再回到原矩阵取相应列；不能直接把 $R$ 的列当作原像空间在基 $F$ 下的坐标。

又因 $Ax=0$ 与 $Rx=0$ 同解，以 $x_4=s$、$x_5=t$ 为自由变量，得到
\[
    x=s\begin{pmatrix}-1\\3\\-2\\1\\0\end{pmatrix}
      +t\begin{pmatrix}3\\-3\\1\\0\\1\end{pmatrix}.
\]
把坐标还原成 $V$ 中的向量，得到核空间的一组基
\[
    -e_1+3e_2-2e_3+e_4,\qquad
    3e_1-3e_2+e_3+e_5.
\]
故 $r(T)=3$、零度为 $2$，二者之和等于 $\dim V=5$。
\end{solve}

一般计算像与核时，可以先将表示矩阵行化简，利用原矩阵的主元列恢复像空间的基，再求 $Ax=0$ 的基础解系并按定义域的基恢复核空间的基。像向量属于陪域、核向量属于定义域，最后一步必须分别使用 $F$ 和 $E$，不能仅给出坐标而省略它们所依赖的基。

\section{不变子空间}
\label{sec:linear-map-invariant-subspace}

把一般线性映射限制到一个子空间后，输出未必还留在该子空间中。对于 $V$ 上的线性变换 $T$，若能找到作用后仍留在其中的子空间，就能在这个较小空间内单独研究 $T$，并通过适当选基使表示矩阵出现分块结构。这是后续研究对角化和相似标准型的起点。

\begin{definition}
\label{def:invariant-subspace}
设 $T\in\mathcal L(V)$，$W$ 是 $V$ 的子空间。若
\[
    T(W)\subseteq W,
\]
则称 $W$ 为 $T$ 的\emph{不变子空间}（invariant subspace）或 $T$-不变子空间。这时 $T|_W:W\to W$ 是 $W$ 上的线性变换，称为 $T$ 在 $W$ 上的限制。
\end{definition}

这里的条件是 $T(W)\subseteq W$，不要求 $T(W)=W$，也不要求 $W$ 中每个向量都被保持不动。$\{0\}$ 和 $V$ 总是不变子空间，称为平凡不变子空间；其余称为非平凡不变子空间。若 $W=\operatorname{span}\{w_1,\ldots,w_r\}$，则检验不变性只需检验每个 $T(w_j)\in W$，因为线性性会使任意线性组合的像也留在 $W$ 中。

\begin{example}
对任意 $T\in\mathcal L(V)$，$\ker T$ 和 $\operatorname{Im}T$ 都是 $T$-不变子空间。若 $T=cI_V$ 是纯量变换，则 $V$ 的每个子空间都是 $T$-不变的。
\end{example}

\begin{solve}
$T(\ker T)=\{0\}\subseteq\ker T$，且 $T(\operatorname{Im}T)\subseteq T(V)=\operatorname{Im}T$。纯量变换则满足 $T(w)=cw\in W$（$w\in W$），所以保持任意子空间。特别地，核中的非零向量可能全被映成零，“不变”并不是逐点固定的意思。
\end{solve}

\begin{example}
实平面上绕原点旋转 $\theta$ 角的变换 $R_\theta$，当 $\theta\notin\pi\mathbb Z$ 时没有非平凡不变子空间；当 $\theta\in\pi\mathbb Z$ 时，每个子空间都不变。
\end{example}

\begin{solve}
实平面的非平凡子空间只能是一条过原点的直线。若 $W=\operatorname{span}\{v\}$（$v\ne0$）不变，则 $R_\theta v=cv$，其中 $c\in\mathbb R$。旋转保持长度，故 $|c|=1$，即向量必须被转到自身或反向，这要求 $\theta\in\pi\mathbb Z$。反过来，这些角对应的变换为 $I$ 或 $-I$，都是纯量变换，故每个子空间都不变。
\end{solve}

\begin{theorem}
\label{thm:invariant-subspace-block-triangular}
设 $\dim V=n$，$W$ 是 $r$ 维子空间，$0<r<n$。取 $W$ 的基 $w_1,\ldots,w_r$，扩充成 $V$ 的基
\[
    E=(w_1,\ldots,w_r,v_{r+1},\ldots,v_n).
\]
则 $W$ 为 $T$-不变子空间，当且仅当 $T$ 在这组基下的矩阵具有分块上三角形式
\begin{equation}
\label{eq:invariant-subspace-block-triangular}
    [T]_E=\begin{pmatrix}A&B\\0&C\end{pmatrix},
\end{equation}
其中 $A$ 为 $r$ 阶方阵，且正是 $T|_W$ 在基 $(w_1,\ldots,w_r)$ 下的表示矩阵。
\end{theorem}

\begin{proof}
若 $W$ 不变，每个 $T(w_j)$ 都是 $w_1,\ldots,w_r$ 的线性组合，其后 $n-r$ 个坐标全为零。因此矩阵前 $r$ 列的下半部分为零，上半部分组成 $A$，其余列组成 $B,C$。反之，若矩阵具有这种形式，则每个 $T(w_j)$ 的后 $n-r$ 个坐标为零，即 $T(w_j)\in W$；由线性性得到 $T(W)\subseteq W$。
\end{proof}

在式~\eqref{eq:invariant-subspace-block-triangular} 中，令 $Z=\operatorname{span}\{v_{r+1},\ldots,v_n\}$，则 $V=W\oplus Z$，但 $Z$ 未必不变。矩阵的右上块 $B$ 正好记录 $Z$ 中向量被映入 $W$ 的分量；$B=0$ 当且仅当这个补空间 $Z$ 也不变。因此有一个不变子空间只保证分块上三角，不能直接断言可以分块对角化。

\begin{corollary}
\label{cor:invariant-direct-sum-block-diagonal}
若有限维空间有直和分解
\[
    V=W_1\oplus\cdots\oplus W_s,
\]
且每个 $W_i$ 都是 $T$-不变子空间，则合并各 $W_i$ 的基得到 $V$ 的一组基，在这组基下
\begin{equation}
\label{eq:invariant-direct-sum-block-diagonal}
    [T]=\operatorname{diag}(A_1,\ldots,A_s),
\end{equation}
其中 $A_i$ 是 $T|_{W_i}$ 的表示矩阵。反之，若矩阵具有这种分块对角形式，相应各组基向量张成的子空间就是不变子空间，并且其直和为 $V$。
\end{corollary}

\begin{proof}
每个 $W_i$ 的基向量的像只含同一个 $W_i$ 中的基向量，故其矩阵列在其他行块中全为零，得到分块对角形式。反之，从各列的零块可知每组基向量的像仍在本组的张成空间内；整组向量是基，保证这些张成空间的和为直和且等于 $V$。
\end{proof}

特别地，$T$ 存在对角表示矩阵，当且仅当 $V$ 是 $n$ 个一维不变子空间的直和。一维子空间 $K v$（$v\ne0$）不变，等价于存在 $\lambda\in K$ 使 $T(v)=\lambda v$；这将引出特征值和特征向量的概念。不过，不变子空间不一定总有不变补空间。例如在基 $(e_1,e_2)$ 下定义 $T(e_1)=0$、$T(e_2)=e_1$，则 $W=K e_1$ 不变。它的任意一维补空间都可写成 $Z=K(ae_1+e_2)$，但 $T(ae_1+e_2)=e_1\notin Z$，所以不存在不变补空间。

\begin{example}
\label{ex:invariant-subspace-textbook}
设 $V$ 的基为 $E=(e_1,e_2,e_3)$，$T$ 在此基下的矩阵为
\[
    A=\begin{pmatrix}
        3&1&-1\\
        2&2&-1\\
        2&2&0
    \end{pmatrix}.
\]
证明 $W=\operatorname{span}\{e_3,e_1+e_2+2e_3\}$ 是 $T$-不变子空间，并求限制变换的表示矩阵。
\end{example}

\begin{solve}
记 $w_1=e_3$、$w_2=e_1+e_2+2e_3$，二者线性无关，构成 $W$ 的一组基。由矩阵各列读出
\[
    T(w_1)=-e_1-e_2=2w_1-w_2,
\]
\[
    T(w_2)=T(e_1)+T(e_2)+2T(e_3)
    =2e_1+2e_2+4e_3=2w_2.
\]
故两个基向量的像都在 $W$ 中，$W$ 不变，并且
\[
    [T|_W]_{(w_1,w_2)}=\begin{pmatrix}2&0\\-1&2\end{pmatrix}.
\]
再以 $e_1$ 扩充为基 $E'=(w_1,w_2,e_1)$，由
\[
    T(e_1)=3e_1+2e_2+2e_3=-2w_1+2w_2+e_1
\]
得到
\[
    [T]_{E'}=\begin{pmatrix}
        2&0&-2\\
        -1&2&2\\
        0&0&1
    \end{pmatrix},
\]
其左下零块正是 $W$ 不变的矩阵表现。
\end{solve}

作为本节结论的直接应用，还可以用不变性的定义检查两个常用性质。若 $W_1,W_2$ 都不变，则 $W_1+W_2$ 与 $W_1\cap W_2$ 也不变：前者由 $T(w_1+w_2)=T(w_1)+T(w_2)$ 得到，后者由交中向量的像同时属于两个子空间得到。若 $S,T$ 满足 $ST=TS$，则 $\ker T$ 和 $\operatorname{Im}T$ 都是 $S$-不变的，因为
\[
    v\in\ker T\ \Longrightarrow\ T(Sv)=S(Tv)=0,
    \qquad S(Tv)=T(Sv)\in\operatorname{Im}T.
\]
这些结论不需要有限维假设；相比之下，若要由 $T$ 可逆且 $W$ 不变推出 $T(W)=W$，就需要额外条件。$W$ 有限维时，$T|_W$ 是单射，由推论~\ref{cor:linear-map-injective-surjective-rank} 也是满射，故 $T(W)=W$，进而 $T^{-1}(W)=W$。






